\title{The Era of 1-bit LLMs: All Large Language Models are in 1.58 Bits}

\begin{abstract}
Breakthroughs like BitNet~\cite{bitnet} are ushering in a new epoch for 1-bit Large Language Models (LLMs). In this study, we propose \textbf{\text{BitNet b1.58}}, a novel 1-bit LLM where each parameter is constrained to a ternary set \{-1, 0, 1\}. Despite this extreme quantization, our model matches the performance of full-precision (FP16 or BF16) Transformer-based LLMs of equivalent scale and training regimen in terms of both perplexity and downstream task accuracy. Moreover, the 1.58-bit formulation delivers pronounced advantages in latency, memory usage, throughput, and energy efficiency. Importantly, it establishes an innovative scaling law and methodology for constructing future LLMs that are simultaneously high-performing and resource-conscious. Additionally, this architecture heralds a shift in computation paradigms, paving the way for specialized hardware tailored for 1-bit LLM processing.
\end{abstract}

\section{The Era of 1-bit LLMs}

The artificial intelligence community has witnessed unprecedented expansion in the scale and performance of Large Language Models (LLMs) in recent years. While these systems have excelled across a broad array of natural language processing applications, their escalating parameter counts present significant deployment obstacles and heighten worries regarding environmental and financial costs due to substantial power usage.

To mitigate these issues, one prevalent strategy involves applying post-training quantization, generating low-bitwidth model variants suitable for inference~\cite{smoothquant, gptq, quip, quip_sharp}. By reducing the numerical precision of both weights and activations, this technique yields significant savings in memory footprint and computational load for LLMs. This trend of quantization has advanced from the standard 16-bit representation toward even more compact forms, including 4-bit models~\cite{gptq, awq}. Nevertheless, post-training quantization, while frequently adopted in production LLM systems, is generally considered suboptimal.

Recent advancements in 1-bit model designs, exemplified by BitNet~\cite{bitnet}, point to a promising path for decreasing the computational demands of large language models (LLMs) without sacrificing accuracy. Traditional LLMs generally operate using 16-bit floating-point numbers (FP16 or BF16), with matrix multiplication constituting the bulk of their computations. As such, the primary computational overhead is incurred through floating-point additions and multiplications. BitNet diverges from this standard by utilizing only integer additions for matrix multiplication, delivering a drastic reduction in energy usage. Since power constraints are often the main bottleneck for chip performance, these energy savings can be leveraged for speedups in computation.

Apart from compute efficiency, model parameter transfer—specifically, moving parameters from off-chip DRAM to the on-chip accelerator memory (like SRAM)—can also pose a significant cost during inference. Although expanding SRAM capacity has been explored to address throughput limitations, this approach can dramatically increase costs relative to DRAM. In comparison, 1-bit LLMs drastically reduce memory requirements, both in terms of storage and bandwidth, relative to their full-precision counterparts. This reduction translates to more efficient and rapid parameter loading from DRAM, enabling swifter and more cost-effective inference.

In this study, we present an advanced 1-bit LLM, \textbf{\text{BitNet b1.58}}, characterized by ternary parameter values drawn from \{-1, 0, 1\}. This is accomplished by supplementing the traditional 1-bit BitNet parameter space with an additional 0 value, thereby resulting in an effective precision of 1.58 bits in binary terms. \text{BitNet b1.58} preserves all of the key benefits introduced by the original 1-bit BitNet, such as its highly efficient computational mechanism, which largely obviates multiplication operations within matrix multiplication and lends itself to substantial optimization. Moreover, its power consumption is on par with the initial BitNet, while surpassing FP16-based LLMs in memory efficiency, throughput, and latency. Additionally, \text{BitNet b1.58} introduces two notable enhancements: first, the added value of 0 in the weight set facilitates explicit feature selection, thus bolstering its modeling strength and potentially yielding notable performance gains over earlier 1-bit LLMs; second, empirical results indicate that, starting from a scale of 3B parameters, \text{BitNet b1.58} achieves comparable perplexity and downstream task results to full-precision (FP16) baselines when controlled for model configuration, size, and training data.

\section{BitNet b1.58}

\text{BitNet b1.58} is developed from the BitNet framework, which implements the Transformer model by substituting the standard \emph{nn.Linear} layer with \emph{BitLinear}. The model is initialized with random parameters and trained with weight precision set to 1.58 bits and activation precision at 8 bits. While rooted in the initial BitNet structure, several revisions have been introduced, which are detailed below.

\paragraph{Quantization Function.} To ensure that the weights are limited to the set \{-1, 0, +1\}, an \emph{absmean} quantization strategy is applied. This approach rescales the weight matrix by dividing by the mean of the absolute values, after which each entry is discretized to the nearest of \{-1, 0, +1\}:

\begin{equation}
    \widetilde{W} = \mathrm{RoundClip}(\frac{W}{\gamma + \epsilon}, -1, 1),
\end{equation}

\begin{equation}
    \mathrm{RoundClip}(x, a, b) = \max(a, \min(b, \mathrm{round}(x))),
\end{equation}

\begin{equation}
    \gamma = \frac{1}{nm}\sum_{ij} |W_{ij}|.
\end{equation}

For activations, quantization follows the procedure outlined in the original BitNet; however, we opt not to rescale activations into the interval $[0, Q_b]$ prior to the application of nonlinear functions. Instead, all activations are normalized per token to the interval $[-Q_b, Q_b]$, effectively bypassing zero-point quantization. This adjustment simplifies both the implementation and optimization at the system level, while empirical evaluations indicate that the resulting performance difference is marginal.

\paragraph{LLaMA-alike Components.}

With the LLaMA~\cite{llama, llama2} architecture widely established as the standard for open-source large language models, \text{BitNet b1.58} is deliberately structured to match this convention and foster compatibility within the open-source community. In particular, our implementation features RMSNorm~\cite{rmsnorm}, SwiGLU~\cite{swiglu}, rotary embedding~\cite{rope}, and omits all biases. As a result, \text{BitNet b1.58} can be incorporated into prevalent open-source toolkits (such as Huggingface, vLLM~\cite{vllm}, and llama.cpp\footnote{\url{https://github.com/ggerganov/llama.cpp}}) with minimal additional engineering.

\section{Results}

We conducted a comparison between \text{BitNet b1.58} and our reproduced FP16 \text{LLaMA LLM} across multiple model scales. Both models were pre-trained on the RedPajama dataset~\cite{redpajama} with 100 billion tokens to maintain experimental parity. Zero-shot evaluations spanned a suite of language understanding benchmarks: ARC-Easy~\cite{arc}, ARC-Challenge~\cite{arc}, Hellaswag~\cite{hellaswag}, Winogrande~\cite{winoGrande}, PIQA~\cite{piqa}, OpenbookQA~\cite{openbookqa}, and BoolQ~\cite{boolq}. Additionally, we assessed validation perplexity on WikiText2~\cite{wikitext2} and C4~\cite{c4} datasets.

For inference efficiency, we measured GPU memory usage and response latency for both \text{LLaMA LLM} and \text{BitNet b1.58}. Evaluations were performed using the FasterTransformer\footnote{\url{https://github.com/NVIDIA/FasterTransformer}} library, which offers high-performance inference for LLMs on GPUs. The \text{BitNet b1.58} runs also leveraged the 2-bit kernel from Ladder~\cite{ladder}. Inference cost was quantified as the generation time per output token, given its centrality to deployment overhead.

As summarized in Table~\ref{tab:ppl}, \text{BitNet b1.58} achieves perplexity comparable to full-precision \text{LLaMA LLM} when scaled to 3B parameters, while reducing inference time by a factor of 2.71 and memory footprint by 3.55. Notably, the 3.9B \text{BitNet b1.58} model operates 2.4 times faster and uses 3.32 times less GPU memory than \text{LLaMA LLM} 3B, while delivering higher performance.

Detailed zero-shot task accuracy can be found in Table~\ref{tab:zero-shot}. Evaluations adhered to the \emph{lm-evaluation-harness}\footnote{\url{https://github.com/EleutherAI/lm-evaluation-harness}} standard workflow. Findings indicate that as the parameter count increases, performance differences between \text{BitNet b1.58} and \text{LLaMA LLM} diminish. Importantly, from the 3B parameter scale, \text{BitNet b1.58} achieves parity with its full-precision counterpart. Consistent with the perplexity outcomes, on downstream tasks \text{BitNet b1.58} 3.9B surpasses \text{LLaMA LLM} 3B, yet requires less computational memory and inference time, marking a clear Pareto advancement relative to established LLM architectures.

\paragraph{Memory and Latency}

We extended our experiments by increasing the model sizes to 7B, 13B, and 70B and analyzed the corresponding computational costs. As presented in Figure~\ref{fig:latency-memory-energy}, both latency and memory requirements exhibit predictable scaling behaviors, with acceleration becoming more pronounced at larger model scales. Notably, at the 70B scale, \text{BitNet b1.58} achieves a 4.1-fold reduction in inference time relative to the \text{LLaMA LLM} reference model. This improvement stems from the scaling of the \emph{nn.Linear} operation’s time complexity with respect to model size. Similarly, the observed pattern for memory usage reflects the fact that the embeddings, which retain full precision, account for a decreasing share of total memory in larger models. Latency and memory measurements were collected with a 2-bit implementation, suggesting additional optimization opportunities remain for further cost reductions.

\paragraph{Energy}

We additionally evaluated the estimated energy expenditure attributable to arithmetic operations in both \text{BitNet b1.58} and \text{LLaMA LLM}. Given that matrix multiplication is the primary source of computational cost in LLMs, our analysis centers on this component. The energy distribution is depicted in Figure~\ref{fig:energy}, showing that most operations in \text{BitNet b1.58} are INT8 additions, whereas \text{LLaMA LLM} relies mainly on FP16 addition and multiplication. Based on the energy model detailed in~\cite{energycost, pokebnn}, the matrix multiplication operations in \text{BitNet b1.58} result in a 71.4-fold reduction in energy consumption on 7nm process hardware. We further provide full-model energy measurements for input sequences of 512 tokens, finding that energy savings grow with increasing model size. As the contribution of the \emph{nn.Linear} component increases at larger scales while the cost of other parts becomes less significant, \text{BitNet b1.58} demonstrates increasingly superior energy efficiency over the FP16-based \text{LLaMA LLM} as models scale up.

\paragraph{Throughput}

To assess throughput, we compare \text{BitNet b1.58} and the \text{LLaMA LLM} with 70B parameters, deploying both on a pair of 80GB A100 GPUs and employing pipeline parallelism~\cite{gpipe} to accommodate the \text{LLaMA LLM} 70B model. The batch size was incrementally increased up to the GPU memory constraints, with a fixed sequence length of 512. As illustrated in Table~\ref{tab:throughput}, \text{BitNet b1.58} 70B is capable of handling batch sizes up to 11 times larger than \text{LLaMA LLM}, which in turn translates into a throughput enhancement by a factor of 8.9.

\textbf{\text{BitNet b1.58} paves the way for a novel scaling law linking model effectiveness and inference resource requirements}. To contextualize, Figure~\ref{fig:latency-memory-energy} and Figure~\ref{fig:energy} provide equivalences for different model sizes under 1.58-bit and 16-bit settings:

\begin{itemize}
    \item The 13B BitNet b1.58 model surpasses a 3B FP16 LLM in latency, memory footprint, and energy efficiency.
    \item The 30B BitNet b1.58 exhibits better latency, memory usage, and energy consumption than a 7B FP16 LLM.
    \item The 70B BitNet b1.58 achieves greater efficiency in all three metrics compared to a 13B FP16 LLM.
\end{itemize}

\paragraph{Training with 2T Tokens}

Token count is a critical element in training LLMs. To evaluate the token scalability of \text{BitNet b1.58}, we conducted training with 2T tokens, utilizing the data protocol from StableLM-3B~\cite{StableLM-3B-4E1T}—the leading open-source model in the 3B parameter class. Performance for both models was measured on a composite benchmark including Winogrande~\cite{winoGrande}, PIQA~\cite{piqa}, SciQ~\cite{sciq}, LAMBADA~\cite{lambada}, and ARC-easy~\cite{arc}. We present zero-shot accuracy results in Table~\ref{tab:2t}; for benchmarks assessed with both accuracy and normalized accuracy, their mean was reported. StableLM 3B's 2T token results are drawn directly from its technical report. Our results demonstrate that \text{BitNet b1.58} consistently outperforms on all evaluated tasks, implying robust generalization in 1.58-bit LLMs as well.

\section{Discussion and Future Work}

\textbf{1-bit Mixture-of-Experts (MoE) LLMs}

Mixture-of-Experts (MoE) architectures have emerged as an efficient way to reduce the computational demands of large language models (LLMs). Although MoE significantly decreases required FLOPs, memory usage and the associated inter-chip data transfers pose significant barriers to practical deployment. The adoption of 1.58-bit LLMs presents a promising solution to these issues. Primarily, compressing the model's memory requirements allows for deploying MoE models on fewer devices. Additionally, the transmission overhead associated with moving activations across chips is also lessened. Ideally, if the reduced footprint enables the model to reside entirely on a single chip, network communication overheads could be eliminated.

\textbf{Native Support of Long Sequence in LLMs}

As LLMs become more prevalent, efficient handling of long sequences has grown increasingly vital. One of the core difficulties in long sequence inference arises from the high memory demands imposed by KV caches. The advent of \text{BitNet b1.58} makes considerable progress in this area by shrinking activation sizes from 16 bits to 8 bits, thereby permitting doubled context lengths on equivalent hardware resources. There is potential to achieve further, lossless compression down to 4 bits or even less with 1.58-bit LLMs—a direction we identify for forthcoming research.

\textbf{LLMs on Edge and Mobile}

Utilizing 1.58-bit LLMs can substantially enhance the effectiveness of language models on edge and mobile hardware, where memory and processing capabilities are generally constrained. The lower memory footprint and energy consumption make it feasible to deploy such models in resource-limited scenarios, unlocking a variety of previously unattainable applications. This not only broadens the functional scope of edge and mobile devices but also introduces fresh use cases for LLM technology. Importantly, 1.58-bit LLMs are better suited for CPU-based computation, which predominates in mobile and edge devices. As a result, \text{BitNet b1.58} can be run efficiently on these platforms, further advancing their potential.

\textbf{New Hardware for 1-bit LLMs}

Initiatives such as Groq\footnote{\url{https://groq.com/}} have demonstrated the advantages and opportunities of building specialized hardware, like LPUs, tailored for LLM inference. Building on this foundation, there is a pressing need to develop novel hardware and system designs specifically tailored to the characteristics of 1-bit LLMs, as enabled by the computation strategies introduced in BitNet~\cite{bitnet}.